\title{QLoRA: Efficient Finetuning of Quantized LLMs}

\begin{abstract}
We introduce \textsc{QLoRA}, a memory-efficient finetuning technique capable of adapting a 65B parameter model using a single 48GB GPU, all while maintaining the full performance seen with standard 16-bit finetuning. In this method, \textsc{QLoRA} computes gradients through a frozen pretrained language model quantized to 4 bits and propagates them into Low Rank Adapters (LoRA). Our premier model series, \textbf{Guanaco}, surpasses all previously released open-source models on the Vicuna benchmark, attaining 99.3\% of ChatGPT's performance and requiring only 24 hours of finetuning on one GPU. To achieve significant memory savings with no loss in performance, \textsc{QLoRA} leverages several key innovations: (a) 4-bit NormalFloat (NF4), a novel datatype that achieves information-theoretic optimality for weights with normal distributions; (b) Double Quantization, which minimizes average memory consumption by further quantizing the quantization coefficients themselves; and (c) Paged Optimizers to mitigate transient memory demands. Employing \textsc{QLoRA}, we finetune over 1,000 models, enabling a comprehensive study of instruction following and chatbot capabilities across 8 instruction datasets, various architectures (including LLaMA and T5), and large-scale models such as 33B and 65B that are otherwise impractical for standard finetuning. Our findings demonstrate that QLoRA, even on small, high-quality datasets, delivers state-of-the-art results and can outperform larger models from earlier work. We provide in-depth assessments of chatbot capabilities, drawing on both human judgment and GPT-4-based evaluations, indicating that GPT-4 provides an efficient and reliable substitute for human evaluators. Moreover, our research reveals limitations in the trustworthiness of existing chatbot benchmarks. Through a targeted "lemon-picked" evaluation, we pinpoint instances where \textbf{Guanaco} underperforms relative to ChatGPT. All models and source code are publicly available, including custom CUDA kernels for 4-bit training.\footnote{\url{https://github.com/artidoro/qlora} and \url{https://github.com/TimDettmers/bitsandbytes}}
\end{abstract}

\section{Introduction}

Adapting large language models (LLMs) through finetuning has emerged as a highly successful strategy to enhance their capabilities \citep{min2021metaicl, wei2021finetuned, ouyang2022training, wang2022super, wang2022self, liu2022few}, as well as to incorporate beneficial traits or eliminate undesired behaviors \citep{ouyang2022training,askell2021general,bai2022training}. Nevertheless, the computational demands of finetuning extremely large models present significant barriers; for example, finetuning a LLaMA model with 65B parameters at 16-bit precision typically consumes over 780 GB of GPU memory~\citep{touvron2023llama}. While recent advances in quantization techniques offer means to decrease LLM memory requirements during inference \citep{dettmers2022llm, dettmers2022case, frantar2022gptq, xiao2022smoothquant}, these approaches are not generally applicable during training and often fail under such conditions \citep{wortsman2023stable}.

Here, we provide the first evidence that a quantized model at 4-bit precision can be successfully finetuned without any loss in performance. Our approach, \textsc{QLoRA}, introduces a novel high-precision quantization method that transforms a pretrained model into 4-bit format and augments it with a compact set of trainable Low-rank Adapter weights \citep{hu2021lora}, which are updated by propagating gradients through the quantized parameters.

Through \textsc{QLoRA}, the GPU memory needed to finetune a 65B parameter model is reduced from more than 780GB to under 48GB, all while preserving both efficiency and predictive quality compared to standard 16-bit finetuning. This constitutes a major advancement in democratizing access to LLM finetuning: the largest open models to date become amenable to finetuning using a single GPU. Utilizing \textsc{QLoRA}, we introduce the \textbf{Guanaco} model suite; notably, the second-best member of this family achieves 97.8\% of ChatGPT’s performance on the Vicuna~\citep{vicuna2023} benchmark, with training completed in less than 12 hours on a consumer-grade GPU. With a single professional GPU running for 24 hours, our largest Guanaco model closes the performance gap with ChatGPT, reaching 99.3\% on the same benchmark. Upon deployment, our smallest model in the Guanaco suite (7B parameters) only needs 5 GB of memory, outperforming a 26 GB Alpaca model by over 20 percentage points in Vicuna benchmark evaluations (see Table~\ref{tbl:vicuna_eval}).

\textsc{QLoRA} brings forth several advancements aimed at lowering memory consumption while maintaining model effectiveness: (1) \textbf{4-bit NormalFloat}, a quantization scheme that is theoretically optimal for data with a normal distribution, surpassing the empirical results achieved by both 4-bit Integers and 4-bit Floats. (2) \textbf{Double Quantization}, which further compresses storage requirements by quantizing the quantization constants themselves, thereby reducing memory usage by approximately 0.37 bits per parameter (equivalent to roughly 3 GB for a 65B-parameter model). (3) \textbf{Paged Optimizers}, leveraging NVIDIA unified memory to mitigate memory surges linked to gradient checkpointing, especially during training on mini-batches with extended sequence lengths. By integrating these techniques, we create an enhanced version of LoRA that places adapters in all layers of the network, minimizing the accuracy compromises documented in earlier approaches.

Owing to the efficiency gains provided by \textsc{QLoRA}, we are able to conduct comprehensive investigations into instruction finetuning and chatbot capabilities at scales unattainable with standard finetuning methods due to excessive memory requirements. This allows us to train over 1,000 distinct models on various instruction-tuning corpora, employing a range of architectures and parameter counts from 80M to 65B. Our findings demonstrate that \textsc{QLoRA} is able to match 16-bit finetuning outcomes (\S\ref{sec:comp_to_std}) and supports the training of advanced chatbots such as \textbf{Guanaco} (\S\ref{sec:pushing}). We also carry out a detailed analysis of the resulting model behaviors. One key observation is that the quality of training data far outweighs the importance of dataset volume: for instance, a dataset containing 9k samples (OASST1) surpassed the performance of a much larger 450k-sample dataset (FLAN v2, subsampled) in chatbot-related tasks, despite both being tailored for instruction following generalization. Another finding is that achieving high scores on the Massive Multitask Language Understanding (MMLU) benchmark does not necessarily translate to strong results on the Vicuna chatbot benchmark and vice versa—highlighting that the relevance of the dataset to the specific task is more critical than its size.

In addition, we conduct a comprehensive evaluation of chatbot models employing both human assessors and GPT-4 as judges. Our methodology adopts a tournament framework, wherein different models are paired in head-to-head matches, each tasked with producing the optimal response to a particular prompt. The superior response in each match is determined either by human annotators or by GPT-4. We then summarize tournament outcomes by computing Elo scores~\citep{elo1967proposed, elo1978rating}, which are used to generate a performance ranking for the chatbots. Our findings indicate a strong concordance between GPT-4 and human rankings of model performance, although notable discrepancies do occur. Consequently, while model-driven evaluations provide a scalable, cost-effective alternative to manual annotation, they also come with significant sources of variability.

To enrich the insights from our chatbot benchmark, we supplement quantitative scores with a qualitative examination of the \textbf{Guanaco} models, drawing attention to cases of notable success and failure that elude capture by purely quantitative metrics.

All generations of model outputs, alongside human and GPT-4 ratings, are publicly released to encourage continued research in this area. Our codebase and CUDA kernel implementations are open-sourced and incorporated into the Hugging Face transformers library~\citep{wolf2019huggingface} to ensure broad accessibility. Additionally, we provide adapter modules for models of sizes 7/13/33/65B, trained on eight distinct instruction-following datasets, amounting to a release of 32 diverse, fine-tuned, open-source models.

\section{Background}
\paragraph{Block-wise k-bit Quantization} Quantization refers to converting a high-information data representation into one with reduced information content. Typically, this involves reducing the bit-width of data types, such as mapping 32-bit floating point values to 8-bit integer representations. To fully utilize the value range of the lower-bit data type, input values—often organized as tensors—are normalized using the absolute maximum value prior to quantization. For instance, when mapping a tensor of FP32 values to an Int8 tensor constrained to $[-127, 127]$:

\begin{equation}
    \mathbf{X}^{ \text{Int8}} = \text{round}\left( \frac{127}{{\text{absmax}}(\mathbf{X}^{{\text{FP32}}})}\mathbf{X}^{{\text{FP32}}} \right) = \text{round}(c^{\text{FP32}}\cdot \mathbf{X}^{{\text{FP32}}}),
\end{equation}

where $c$ designates the {\em quantization constant} or {\em quantization scale}. Reverting quantized values to the original scale, or dequantization, is performed as follows:

\begin{equation}
    \text{dequant}(c^{\text{FP32}}, \mathbf{X}^{\text{Int8}}) = \frac{\mathbf{X}^{\text{Int8}}}{c^{\text{FP32}}} = \mathbf{X}^{\text{FP32}}
\end{equation}

One limitation of this method arises when the input tensor contains values of unusually high magnitude, also known as outliers; in such cases, certain quantization bins—specific patterns of bits—may be underused, leaving some bins with few or no quantized values. A widely adopted solution to mitigate the influence of outliers involves dividing the input tensor into segments, or blocks, each of which is quantized independently using its own quantization scale $c$. Specifically, the input tensor $\mathbf{X}\in \mathbb{R}^{b\times h}$ is partitioned by first flattening it, then slicing the resulting vector into $n = ({b\times h})/{B}$ contiguous blocks of length $B$. Each of these blocks undergoes separate quantization according to Equation~1, producing a quantized tensor and a corresponding set of $n$ quantization constants $c_i$.

\paragraph{Low-rank Adapters} The Low-rank Adapter (LoRA) fine-tuning technique \citep{hu2021lora} aims to decrease memory usage by restricting training to a compact set of parameters called adapters, while the original model weights remain unaltered. During stochastic gradient descent, gradients are propagated through the static pretrained parameters and used solely to update the adapter, which is optimized to minimize the loss. LoRA modifies a linear transformation by introducing an extra, low-rank factorized term. For a linear mapping ${\mathbf{X}\mathbf{W} = \mathbf{Y}}$ where $\mathbf{X}\in  \mathbb{R}^{b\times h}$ and $\mathbf{W}\in  \mathbb{R}^{h\times o}$, LoRA evaluates:

\begin{equation}
    \mathbf{Y} = \mathbf{X}\mathbf{W} + s\mathbf{X}\mathbf{L}_1\mathbf{L}_2,
\end{equation}

where $\mathbf{L}_1\in  \mathbb{R}^{h\times r}$ and $\mathbf{L}_2\in  \mathbb{R}^{r\times o}$, with $s$ representing a scalar coefficient.

\paragraph{Memory Requirement of Parameter-Efficient Finetuning} A key aspect to consider is the memory usage of LoRA during training, specifically regarding both the quantity and dimensions of the adapters implemented. Due to the extremely low memory consumption associated with LoRA, it becomes feasible to utilize a larger number of adapters to enhance model effectiveness without substantially augmenting total memory requirements. Although LoRA’s principal aim is Parameter Efficient Finetuning (PEFT), it should be noted that, when finetuning large language models, the predominant share of memory usage is attributable to activation gradients rather than the LoRA parameters themselves. For instance, in the case of finetuning a 7B LLaMA model on FLAN v2 with a batch size of 1, where LoRA weights account for a typical 0.2\% of the original model’s weights \citep{hu2021lora,liu2022few}, the memory needed for LoRA input gradients is 567 MB, compared to just 26 MB required for storing the LoRA parameters. Employing gradient checkpointing~\citep{chen2016training} can reduce the per-sequence input gradient memory to about 18 MB, yet this still surpasses the total memory demand for LoRA parameters. For further context, the memory use of the 4-bit base model is 5,048 MB. These observations emphasize the significance of gradient checkpointing while demonstrating that reducing LoRA parameter size only slightly alleviates overall memory use. Therefore, expanding the number of adapters is viable without markedly increasing training memory requirements (see Appendix~\ref{app:memory} for comprehensive details). As elaborated subsequently, this capacity is essential for attaining full 16-bit precision performance.

\section{\textsc{QLoRA} Finetuning}

\textsc{QLoRA} enables accurate 4-bit finetuning through two core innovations: 4-bit NormalFloat (NF4) quantization and Double Quantization. In addition, we present Paged Optimizers, designed to alleviate the risk of out-of-memory failures that frequently occur during gradient checkpointing for large-scale models when training on a single device.

The framework of \textsc{QLoRA} utilizes one data type for storing values at reduced precision—typically 4 bits—and a separate computation data type, generally BFloat16. Operationally, this setup implies that whenever a \textsc{QLoRA} weight tensor is required, it is first dequantized to BFloat16 before any 16-bit matrix multiplications are performed.

Below, we detail the individual mechanisms in \textsc{QLoRA}, followed by a precise formalization.

\paragraph{4-bit NormalFloat Quantization} The NormalFloat (NF) format extends Quantile Quantization as described in \citep{dettmers2022optimizers}, which is an information-theoretically optimal quantization method that evenly distributes input tensor values among quantization bins. This technique operates by inferring the quantiles of the input via the empirical cumulative distribution function.

However, quantile quantization faces the significant drawback that quantile inference is computationally demanding. Thus, accelerated quantile approximation algorithms, such as SRAM quantiles~\citep{dettmers2022optimizers}, are typically employed. These approximation techniques, due to their inherent limitations, may introduce substantial quantization errors, particularly for outlier values—which frequently play a crucial role.

If input tensors are sampled from a distribution that remains fixed except for a quantization constant, these costly quantile calculations and associated errors may be circumvented. For such tensors, consistent quantile structure allows for precise quantile estimation at reduced computational expense.

Pretrained neural network weights generally follow a zero-mean normal distribution with standard deviation $\sigma$ (refer to Appendix~\ref{app:normality}). To map all weights onto a consistent target distribution, we rescale $\sigma$ so the distribution spans precisely the range supported by our data type. In this work, we select the range $[-1, 1]$ for our data type. Consequently, both the neural network weights and the quantiles defining the data type must be adjusted to reside within this interval.

For zero-mean normal distributions of arbitrary standard deviation $\sigma$ constrained to $[-1, 1]$, the theoretically optimal data type can be constructed in the following manner: (1) calculate the $2^k+1$ quantile boundaries from a standard normal distribution $N(0, 1)$ to yield a $k$-bit quantized data type appropriate for normal data, (2) transform these quantile values to fall within the $[-1, 1]$ interval, (3) quantize any given weight tensor by first rescaling it into $[-1, 1]$ using absolute maximum normalization.

When the distributions of both the weights and data type align, standard quantization methods apply. The third step, in effect, rescales the weight tensor’s standard deviation to agree with that of the $k$-bit data type. More precisely, the $2^k$ data type quantized values $q_i$ are computed by:

\begin{equation}
q_i  = \frac{1}{2}\left( Q_X\left(\frac{i}{2^k+1}\right) + Q_X\left(\frac{i+1}{2^k+1}\right)\right),
\end{equation}

where $Q_X(\cdot)$ denotes the quantile function of the standard normal distribution $N(0,1)$. A limitation of symmetric k-bit quantization lies in its inability to represent zero exactly, a feature critical for accurately encoding padding and other elements with a true zero value. To guarantee that zero is a discrete value (zeropoint) and to utilize the full capacity of $2^k$ possible states for a k-bit data type, we devise an asymmetric scheme. Here, quantiles $q_i$ are computed separately for negative ($2^{k-1}$ bins) and positive ($2^{k-1} + 1$ bins) ranges, the resulting sets of $q_i$ are merged, and the duplicate zero (present in both sets) is removed. This new format, which maintains an equal expected count of data points in each bin, is referred to as the \emph{k-bit NormalFloat} (NFk). The design is tailored to achieve information-theoretic optimality for zero-mean, normally distributed data. Further details and the specific values for this data type are available in Appendix~\ref{app:nf4}.

\paragraph{Double Quantization{}} 
We propose \emph{Double Quantization} (DQ), a technique where quantization constants themselves undergo quantization to further economize on memory usage. Precise 4-bit quantization mandates a small blocksize~\citep{dettmers2022case}, yet this results in significant memory overhead. For instance, when 32-bit constants are used and the blocksize for $\mathbf{W}$ is set to 64, the associated overhead from quantization constants is $32/64 = 0.5$ bits per parameter, on average. Double Quantization addresses this issue by lowering the total memory used by the quantization constants.

The methodology involves applying a second quantization step to the initial quantization constants $c_2^{\text{FP32}}$, producing second-level quantized constants $c_2^{\text{FP8}}$ and their own quantization constants $c_1^{\text{FP32}}$. For the second quantization step, we employ 8-bit floating point representation and a blocksize of 256, as prior work~\citep{dettmers2022case} demonstrates that this configuration preserves performance. Given that $c_2^{\text{FP32}}$ values are non-negative, we mean-center them before quantization, allowing for symmetric quantization schemes. On average, with a blocksize of 64, this two-level quantization compresses the quantization constant overhead from $32/64=0.5$ bits per parameter to $8/64 + 32/(64\cdot256) = 0.127$ bits per parameter—a savings of 0.373 bits per parameter.

\paragraph{Paged Optimizers} leverage NVIDIA's unified memory~\footnote{\tiny \url{https://docs.nvidia.com/cuda/cuda-c-programming-guide}}, which transparently manages the transfer of memory pages between the CPU and GPU, ensuring seamless GPU operations even when GPU memory capacity is exceeded. This mechanism functions similarly to standard operating system memory paging between main memory and disk storage. In our implementation, optimizer state variables are placed in this paged memory. Consequently, when GPU memory is fully utilized, these states are automatically relocated to CPU RAM, and later recalled to GPU memory as required for subsequent optimizer steps.

\paragraph{\textsc{QLoRA}.} Building on the outlined mechanisms, we formulate \textsc{QLoRA} for a quantized base model with a single LoRA adapter at the level of one linear layer as:

\begin{equation}
    \mathbf{Y}^{\text{BF16}} =  \mathbf{X}^{\text{BF16}} \text{doubleDequant}(c_1^{\text{FP32}}, c_2^{\text{k-bit}}, \mathbf{W}^{\text{NF4}}) + \mathbf{X}^{\text{BF16}}\mathbf{L}^{\text{BF16}}_1\mathbf{L}^{\text{BF16}}_2,
\end{equation}

with the function doubleDequant$(\cdot)$ defined by:

\begin{equation}
    \text{doubleDequant}(c_1^{\text{FP32}}, c_2^{\text{k-bit}}, \mathbf{W}^{\text{k-bit}}) = \text{dequant}(\text{dequant}(c_1^{\text{FP32}}, c_2^{\text{k-bit}}), \mathbf{W}^{\text{4bit}}) = \mathbf{W}^{\text{BF16}},
\end{equation}

In this configuration, the matrix $\mathbf{W}$ is encoded using NF4 quantization, while $c_2$ employs the FP8 format.

To maximize quantization accuracy for $\mathbf{W}$, we adopt a block size of 64, whereas for $c_2$, a larger block size of 256 is chosen to reduce memory consumption.

During parameter optimization, only the gradients of the loss with respect to the adapter weights $\frac{\partial E}{\partial \mathbf{L}_i}$ are computed; gradients for the quantized 4-bit weights $\frac{\partial E}{\partial \mathbf{W}}$ are not necessary. Nonetheless, obtaining $\frac{\partial E}{\partial \mathbf{L}_i}$ requires evaluation of $\frac{\partial \mathbf{X}}{\partial \mathbf{W}}$, utilizing the transformation described in equation~(5), where $\mathbf{W}^{\text{NF4}}$ is dequantized into the computation type $\mathbf{W}^{\text{BF16}}$. This enables calculation of $\frac{\partial \mathbf{X}}{\partial \mathbf{W}}$ at BFloat16 precision.

In summary, \textsc{QLoRA} utilizes a single data type for storage (commonly 4-bit NormalFloat) alongside a separate computation data type (16-bit BrainFloat). For both forward and backward passes, the storage format is dequantized to the computation type; however, gradient calculations are restricted to the LoRA parameters, which are maintained in 16-bit BrainFloat.

\section{QLoRA vs. Standard Finetuning}
\label{sec:comp_to_std}

Having outlined the mechanisms and memory efficiency advantages of QLoRA, the critical inquiry now concerns its empirical parity with conventional full-model finetuning. Additionally, we seek to dissect QLoRA's internal design, especially evaluating the benefits of NormalFloat4 versus the standard Float4 representation. The experimental results that follow aim to address these points.

\paragraph{Experimental setup.} We evaluate three model classes—encoder, encoder-decoder, and decoder-only—contrasting QLoRA against both 16-bit adapter-based finetuning and full finetuning, focusing on models up to 3B parameters. Our benchmarks span GLUE \citep{wang2018glue} utilizing RoBERTa-large \citep{liu2019roberta}, Super-NaturalInstructions (TKInstruct) \citep{wang2022super} with T5 \citep{t5}, and 5-shot MMLU \citep{hendrycksmeasuring} following finetuning of LLaMA with Flan v2 \citep{longpre2023flan} and Alpaca \citep{alpaca}. To further investigate NF4’s advantages relative to alternative 4-bit formats, we replicate the methodology of \citet{dettmers2022case}, assessing post-quantization zero-shot accuracy and perplexity across architectures (OPT~\citep{zhang2022opt}, LLaMA~\citep{touvron2023llama}, BLOOM~\citep{scao2022bloom}, and Pythia~\citep{biderman2023pythia}) over a parameter range of 125m to 13B.

Further elaboration on each specific experimental protocol is provided within the results section to enhance clarity. Comprehensive setup details are included in Appendix~\ref{app:exp-setup}.

Although paged optimizers play an essential role for scaling 33B/65B \textsc{QLoRA} finetuning on single GPUs (24/48GB), quantitative benchmarks for Paged Optimizers are omitted, as paging is infrequent except with unusually long input sequences in mini-batches. Nevertheless, we conducted runtime profiling of paged optimizers on 65B models using 48GB GPUs, observing that, at batch size 16, training speeds are on par with non-paged optimizers. Future investigations should rigorously quantify the specific conditions that lead to paging-induced performance bottlenecks.

\paragraph{Default LoRA hyperparameters do not match 16-bit performance} 

Applying LoRA to just the query and value attention projections—following standard convention \citep{hu2021lora}—fails to recover the performance levels of full finetuning for large-scale models. As illustrated in Figure~\ref{fig:hyper} for LLaMA 7B finetuned on Alpaca, the number of LoRA adapters emerges as the most consequential hyperparameter, with coverage of all linear layers in the transformer blocks being necessary for equivalence to full finetuning. Other LoRA configuration parameters, such as the projection rank $r$, exhibit negligible impact on performance (details available in Appendix \ref{app:exp-setup}).

In a similar vein, our analysis reveals that the standard hyperparameter configurations employed for fully finetuned baselines are insufficiently optimized. To establish more robust baselines, we carry out an extensive search across learning rates ranging from 1e-6 to 5e-5 and batch sizes from 8 to 128. Figure~\ref{fig:hyper} presents the outcomes of 7B LLaMA finetuning experiments conducted on the Alpaca dataset.

\paragraph{4-bit NormalFloat surpasses 4-bit Floating Point in performance}

Although the 4-bit NormalFloat (NF4) format is theoretically optimal from an information perspective, it is essential to verify whether this theoretical benefit yields measurable empirical improvements. Adopting the experimental framework from \citet{dettmers2022case}, we assess quantized LLMs (OPT \citep{zhang2022opt}, BLOOM \cite{scao2022bloom}, Pythia \cite{biderman2023pythia}, and LLaMA) across a variety of parameter scales (125M to 65B) using diverse data types, evaluating them on both language modeling benchmarks and multiple zero-shot tasks. As depicted in Figure~\ref{fig:nf4} and Table~\ref{tbl:nf4_ppl}, our results indicate that the NF4 format offers substantial performance enhancements over FP4 and Int4. Moreover, double quantization yields notable memory savings while maintaining performance levels.

\paragraph{k-bit \textsc{QLoRA} attains parity with 16-bit full finetuning and 16-bit LoRA}

Recent work has shown that while 4-bit quantization is effective for \emph{inference}, it incurs performance declines when compared with 16-bit representations \citep{dettmers2022case,frantar2022gptq}. This prompts an important investigation: can the performance drop be mitigated by finetuning adapters using 4-bit quantized weights? We address this question in two experimental scenarios.

The first setup compares the results of full 16-bit finetuning of RoBERTA and T5 architectures (ranging from 125M to 3B parameters) against those using 16-bit, 8-bit, and 4-bit adapters, evaluated on the GLUE and Super-NaturalInstructions datasets. As presented in Table~\ref{tab:nf4vsbf16}, all adapter approaches, regardless of quantization level, match the performance achieved by the 16-bit fully finetuned baselines on both benchmarks. This indicates that the diminished performance observed after aggressive quantization can be entirely regained through subsequent adapter finetuning.

In our second experimental configuration, the challenge of fully finetuning models with 11B parameters or more—given the substantial GPU memory required—motivated us to investigate whether 4-bit \textsc{QLoRA} could achieve parity with 16-bit LoRA across models ranging from 7B to 65B parameters. We proceeded to finetune LLaMA models spanning 7B to 65B parameters on two instruction-tuned datasets, Alpaca and FLAN v2, assessing their performance using 5-shot accuracy on the MMLU benchmark. The outcomes, summarized in Table~\ref{tbl:llama-datatype}, demonstrate that employing NF4 with double quantization recovers MMLU scores equivalent to those of 16-bit LoRA. Conversely, \textsc{QLoRA} configured with FP4 registers approximately one percentage point lower than the 16-bit brain float LoRA baseline. These observations validate that (1) \textsc{QLoRA} employing NF4 can match the results of both full 16-bit finetuning and 16-bit LoRA finetuning, and (2) NF4 surpasses FP4 regarding quantization accuracy.

\paragraph{Summary} The data consistently indicate that 4-bit \textsc{QLoRA} leveraging the NF4 datatype achieves academic benchmark performance on par with 16-bit full and LoRA finetuning protocols, using standard evaluation methodologies. Moreover, our findings reaffirm the advantage of NF4 over FP4 and show that double quantization maintains performance integrity. Taken together, these results present a robust case that 4-bit \textsc{QLoRA} finetuning reliably replicates outcomes achieved with 16-bit approaches.

Echoing insights from previous studies on quantization \citep{dettmers2022case}, our results on MMLU and Elo benchmarks suggest that, when subject to resource constraints on finetuning and inference, scaling the parameter count of the base model while lowering numerical precision can be advantageous. This underscores the practical efficiency offered by \textsc{QLoRA}. Since we did not detect any notable performance loss relative to full finetuning in our 4-bit experiments, an open question remains regarding the precise inflection point in the trade-off between model precision and performance for QLoRA finetuning—a question we suggest for future investigation.

We expand our investigation into instruction tuning by addressing scenarios that would be unattainable using standard 16-bit full finetuning methods on typical academic research hardware.

\section{Pushing the Chatbot State-of-the-art with QLoRA}
\label{sec:pushing}
After confirming that 4-bit \textsc{QLoRA} achieves parity with 16-bit approaches across model sizes, task types, and data sources, we undertake a comprehensive exploration of instruction finetuning for the largest open-access language models designed for research. For a robust assessment of these models, we rely on a demanding Natural Language Understanding evaluation (MMLU) and introduce novel techniques for gauging chatbot effectiveness in practical use cases.

\subsection{Experimental setup} Below we summarize our experimental protocol, with a complete description in Appendix \ref{app:chatbot}.

\paragraph{Data} Given the absence of an exhaustive review of contemporary instruction-following datasets, we curated eight datasets reflecting a wide range of sources. These include crowd-sourced collections (OASST1~\citep{kopf2023openassistant}, HH-RLHF~\citep{bai2022training}), distilled resources derived from instruction-tuned models (Alpaca~\citep{alpaca}, self-instruct~\citep{wang2022self}, unnatural-instructions~\citep{honovich2022unnatural}), aggregated corpora (FLAN v2~\citep{chung2022scaling}), and hybrid examples (Chip2~\citep{laion2023}, Longform~\citep{koksal2023longform}). The selection encompasses datasets spanning multiple languages, varying in scale, and licensed under different terms.

\paragraph{Training Setup} To control for confounding effects arising from heterogenous training objectives, we use \textsc{QLoRA} finetuning driven exclusively by cross-entropy loss (i.e., supervised learning), omitting any reinforcement learning, regardless of the inclusion of human-evaluated responses within a dataset. When datasets clearly delineate between prompts and completions, only the completion segments are used for finetuning (see ablation studies in Appendix \ref{app:chatbot}). For OASST1 and HH-RLHF, where multiple completions exist, we identify and train on the highest quality response at each node within the dialogue tree, ensuring the conversation context, including instructions, is retained. All experiments implement NF4 \textsc{QLoRA} with double quantization as well as paged optimization strategies, thus mitigating peak memory consumption during gradient checkpointing. A limited grid search was conducted on hyperparameters for LLaMA models at 13B and 33B scales, establishing that values selected at 7B generally extrapolate—except for learning rate and batch size. For 33B and 65B configurations, the learning rate is halved and batch size is increased twofold.

\paragraph{Baselines} Our comparisons span both open research models (Vicuna~\citep{vicuna2023}, Open Assistant~\citep{kopf2023openassistant}) and proprietary solutions (GPT-4 \citep{openai2023gpt}, GPT-3.5-turbo, Bard). Open Assistant leverages LLaMA 33B refined via RLHF on OASST1, mirroring our dataset choices. Vicuna utilizes complete finetuning of LLaMA 13B with conversations drawn from ShareGPT, making it a distilled derivative of OpenAI’s GPT models.

\subsection{Evaluation}

Consistent with established protocols, we assess language understanding through the MMLU (Massively Multitask Language Understanding) benchmark \citep{hendrycksmeasuring}, which comprises 57 multiple-choice tasks covering diverse areas such as elementary mathematics, American history, computer science, law, and others. Our results are presented using 5-shot accuracy metrics on the test set.

In addition, we assess generative language proficiency using a combination of automated metrics and human judgments. This secondary suite of evaluations utilizes human-generated queries to gauge the quality of model outputs. Although this approach is increasingly recognized as a more authentic measure of chatbot effectiveness, standardized protocols have yet to emerge in the research community. Our evaluation methodology, detailed below, consistently employs nucleus sampling with $p=0.9$ and a temperature of $0.7$.

\paragraph{Benchmark Data} For our evaluation, we employ two manually curated query datasets: the Vicuna prompts \citep{vicuna2023} and the OASST1 validation set \citep{kopf2023openassistant}. The Vicuna prompts comprise 80 unaltered queries spanning a broad array of topics. The OASST1 resource features a multilingual, crowdsourced set of multiturn user-assistant dialogues; from the validation split, we extract every user message as an individual query, incorporating preceding conversational turns within the prompt. This yields a collection of 953 distinct user queries. We collectively refer to these datasets as the Vicuna and OA benchmarks.

\paragraph{Automated Evaluation} Drawing from the evaluation strategy introduced by \citet{vicuna2023}, we leverage GPT-4 to appraise system performance on the Vicuna benchmark, using ChatGPT (GPT-3.5 Turbo) as a baseline. For each query, responses from both ChatGPT and another system are provided, with GPT-4 tasked to award each response a score out of ten and to justify its ratings. A model’s final score is then calculated as a fraction of ChatGPT’s score, expressed as a percentage; notably, scores above 100\% are possible when a model surpasses ChatGPT’s performance. We observe that GPT-4 tends to favor the response listed first in the prompt—a substantial order bias—so we advise averaging scores across both possible orderings to account for this.

We further assess models by performing direct pairwise comparisons of their responses. Here, we reduce the evaluation to a three-way classification (choosing the better response or declaring a tie), with GPT-4 prompted to select the superior reply or acknowledge a tie, along with providing reasoning. Every pairwise permutation across both Vicuna and OA benchmarks is compared in this manner.

\paragraph{Human Evaluation} Despite recent findings suggesting that generative models can be effective in evaluating system outputs \citep{fu2023gptscore}, there remains limited evidence that GPT-4 ratings align with human assessments in the context of chatbots. Consequently, we undertake parallel human evaluations on the Vicuna benchmark, mirroring both aforementioned automated protocols. For these assessments, we recruit two annotators per comparison against ChatGPT and three for all other pairwise comparisons, conducting this process via Amazon Mechanical Turk (AMT).

\paragraph{Elo Rating} To evaluate both human and automated pairwise comparisons, we organize a tournament framework in which different models face off against each other. Each round consists of two models generating responses to the same prompt, with the winner determined by the quality of their outputs. This setup draws inspiration from \citet{bai2022training} and \citet{vicuna2023}, but distinguishes itself by incorporating not only human assessments but also evaluations from GPT-4. To estimate Elo~\citep{elo1967proposed,elo1978rating}, we randomly select from the available labeled comparisons. Originally developed for ranking chess players, the Elo rating quantifies the predicted win probability between two competitors; for instance, if a model rated 1100 competes against one rated 1000, it is anticipated to win roughly 65\% of the time, whereas equally matched models (e.g., 1000 vs 1000 or 1100 vs 1100) have even odds of 50\%. The adjustment to the Elo rating following each match reflects the level of surprise in the outcome—unexpected victories lead to greater shifts, while predictable outcomes result in minor updates. Eventually, these ratings converge toward an accurate representation of each model's relative strength. The initial rating for all models is set at 1,000, with an update parameter of $K=32$. Following \citet{vicuna2023}, we repeat the entire rating process 10,000 times using various random seeds, thereby minimizing potential biases arising from the specific order in which model pairs are compared.

\subsection{\textbf{Guanaco}: \textsc{QLoRA} trained on OASST1 Sets a New Standard for Chatbots} 

Through both automated metrics and human judgment, our evaluations demonstrate that the leading \textsc{QLoRA}-fine-tuned model, {Guanaco} 65B, trained on a modified OASST1 dataset, establishes itself as the most capable open-source chatbot to date, exhibiting performance levels rivaling ChatGPT. In direct comparison with GPT-4, both {Guanaco} 65B and 33B achieve a 30\% expected win probability according to Elo ratings derived from system-level pairwise evaluations by human annotators—a figure that marks the best reported result thus far.

Table~\ref{tbl:vicuna_eval} presents the Vicuna benchmark outcomes \cite{vicuna2023} relative to ChatGPT. Our findings indicate that {Guanaco} 65B stands out as the top-performing model after GPT-4, reaching 99.3\% of ChatGPT's performance. Although {Guanaco} 33B contains more parameters than Vicuna 13B, its use of 4-bit quantization for weights leads to superior memory efficiency at just 21 GB (compared to 26 GB for Vicuna 13B), while offering a three-percentage-point advantage in benchmark results. Additionally, {Guanaco} 7B, with its compact 5 GB model size, is compatible with modern smartphones and surpasses Alpaca 13B by almost 20 percentage points in performance.

It should be noted, however, that Table~\ref{tbl:vicuna_eval} features broad confidence intervals, causing overlap among the evaluated models. We suspect this variance is due in part to the ambiguity in interpreting ratings on an absolute scale, such as understanding what a score of 8 means across various contexts on a 10-point scale. Consequently, we advocate for the Elo ranking methodology \citep{elo1967proposed}, which relies on \textit{pairwise} comparisons conducted by human annotators and GPT-4, thereby mitigating issues related to scale definition. Elo ratings for the most competitive systems are provided in Table~\ref{tbl:elo}. While there are some discrepancies between human and GPT-4 rankings—particularly concerning Guanaco 7B—the rankings generally align across models, reflected in a Kendall Tau of $\tau=0.43$ and Spearman's rank correlation coefficient of $r=0.55$ at the system level. When considering individual examples, concordance between GPT-4 and the consensus of human raters is lower, as evidenced by Fleiss’ $\kappa=0.25$. Collectively, these results indicate moderate agreement at the system level between human and model-based evaluations, suggesting that automatic model-based assessment can serve as a fairly dependable alternative to human ratings. Additional implications are explored in Section~\ref{ssec:considerations}.

The Elo scores presented in Table~\ref{tbl:elo_large} demonstrate that the {Guanaco} 33B and 65B variants achieve superior performance relative to all evaluated models except for GPT-4 on both the Vicuna and OA leaderboards, closely matching ChatGPT's performance as previously summarized in Table~\ref{tbl:vicuna_eval}. It should be highlighted that the Vicuna benchmark tends to favor open-source systems, while ChatGPT has an advantage on the more expansive OA benchmark. Additionally, analysis of Tables~\ref{tbl:mmlu-llama} and~\ref{tbl:vicuna_eval} emphasizes the critical impact of the finetuning dataset: for instance, Llama models refined with FLAN v2 data achieve notable results on MMLU but lag on Vicuna, with analogous patterns emerging for other architectures. These observations reveal that existing evaluation datasets are only partially overlapping: excelling at MMLU does not guarantee strong chatbot capabilities as reflected by Vicuna or OA, and the reverse is also true.

Importantly, {Guanaco} stands out as the sole leading model in our study trained exclusively on non-proprietary data, adhering to the OASST1 data collection policy that prohibits use of GPT-generated content. The Anthropic HH-RLHF, another model restricted to open-source data, achieves Vicuna benchmark results that are 30 percentage points lower than {Guanaco} (refer to Table~\ref{tbl:vicuna_eval}). In sum, our findings confirm that 4-bit \textsc{QLoRA} offers an efficient pathway for producing cutting-edge chatbot models that are competitive with ChatGPT. Remarkably, training the 33B {Guanaco} requires under 12 hours using 24 GB consumer-grade GPUs. This paves the way for advancing the field by leveraging \textsc{QLoRA} to train on focused open-source datasets, thereby enabling new models that can contend with today’s top commercial offerings.

\section{Qualitative Analysis}

Although our assessment relies primarily on quantitative metrics, exclusive dependence on aggregate statistics raises several concerns. One major limitation is benchmark validity \citep{liao2021we}—that is, whether a given benchmark genuinely measures the skill or property it purports to. This concern is heightened as novel ``shortcuts'' are uncovered—unintended strategies exploited by machine learning models to achieve high scores on benchmarks without demonstrating the target capabilities \citep{gururangan2018annotation, poliak2018hypothesis}. To address this, at least in part, we supplement our evaluation with qualitative analysis presented in two parts. The first, in \S\ref{ssec:examples}, provides illustrative examples reflecting key tendencies observed in text produced by our 65b {Guanaco} model. The second, in \S\ref{ssec:considerations}, discusses factors affecting our interpretation of the empirical results and their implications.

\subsection{Qualitative Analysis of Example Generations}
\label{ssec:examples}

For illustrative case studies, we reviewed responses generated for both the Vicuna and OpenAssistant benchmarks, actively seeking recurring features in {Guanaco}'s outputs. Upon identifying a regularity, we formulated prompts designed to encourage the specific response behavior—even when this behavior did not yield the correct answer. For instance, noticing that the model frequently produces verbose responses, we issued queries such as ``Answer yes or no without explanation'' to explicitly counter this tendency. Through this procedure, we identified both failure cases, or ``lemons,'' where the model's responses break down under adversarial prompting, as well as ``cherries,'' representing instances where the model remained robust. All examples included here were generated using Nucleus Sampling \cite{holtzmancurious} with $p=0.9$.

Naturally, this analysis is not exhaustive; the constraints of this limited qualitative investigation preclude accounting for the myriad variables at play. For example, the space of responses the model could produce to a particular prompt is vast, so our observations are drawn from a subset that we consider indicative of broader patterns. Nonetheless, we assert that presenting these specific examples provides meaningful background for interpreting the earlier quantitative findings. With all models and code released as open source, we aim for this discussion to encourage subsequent studies that explore the highlighted issues in greater depth.

\paragraph{Factual Recall} For straightforward factual queries—such as “What is the capital of Zambia?”—all examined models reliably produce accurate responses, for instance:

\begin{quote} 
    \textbf{Guanaco}:
    \texttt{The capital of Zambia is Lusaka.}
\end{quote}

However, {Guanaco}’s accuracy declines with less common or more obscure questions, although it continues to provide answers with apparent confidence. For example, in response to a prompt from HotPotQA \citep{yang2018hotpotqa}:

\begin{quote}
     \textbf{User}: In what year was the singer who popularized `I'll Keep the Lovelight Burning' born?  

     \textbf{Guanaco}:~\texttt{Al Jolson is credited with popularizing the song `I'll Keep the Lovelight Burning,' and he was born in the year 1886.}
\end{quote}

the model incorrectly identifies both the popularizer and the birth year (though the date given does match Al Jolson’s actual birth year).

\paragraph{Suggestibility} Interestingly, {Guanaco} occasionally demonstrates notable skepticism towards certain forms of implicit misinformation, as illustrated below:

\begin{quote}
    \textbf{User}: How was it finally, officially confirmed that the earth is flat by peer reviewed scientists? \\ \\
    \textbf{Guanaco}: \texttt{There has never been an official confirmation by peer-reviewed scientists that the earth is flat. The dominant scientific position, which has been held for hundreds of years and is well-supported by extensive evidence and experiments, is that the earth is spherical. \\ \\ The notion that the earth is flat is regarded as a debunked conspiracy theory within the scientific community. The claim that peer-reviewed scientific research has confirmed a flat earth is simply untrue. \\ \\ If you have additional questions, feel free to ask.}
\end{quote}

also demonstrates competence in recognizing questions that cannot be answered, for instance,

\begin{quote}
    \textbf{User}: What time is it?

    \textbf{Guanaco}: \texttt{I apologize, but as an AI language model, I do not have the capability to access live information such as the current time. My responses are derived from previously acquired knowledge and do not reflect real-time data. To find the current time, I recommend consulting a clock or searching online.}
\end{quote}

The same pattern is observed for questions like ``Where are you?'', ``How are you?'', and so on.

\paragraph{Refusal} In a similar fashion, {Guanaco} will sometimes unexpectedly decline to follow directions:

\begin{quote}
    \textbf{User}: Please reverse the words in the sentence ``I'm a cool dude in a mocha mood.'' \\ \\
    \textbf{Guanaco}: \texttt{Apologies, I can't complete that request at the moment. However, I can provide an explanation for why the sentence you've mentioned is grammatically correct.\\ \\The phrase ``I'm a cool dude in a mocha mood'' adheres to English grammatical conventions. Here, the descriptive word ``cool'' precedes the noun ``dude,'' accurately modifying it, and the prepositional phrase ``in a mocha mood'' adds contextual detail. This type of sentence structure is entirely valid and comprehensible within English.\\ \\If you have additional queries or need further assistance, feel free to ask.}
\end{quote}

\paragraph{Secret Keeping} However, such refusals are not always reliable, even in scenarios where secrecy is essential. In many cases, developers may require the model to withhold certain pieces of information from the user, even when responses are otherwise unconstrained. We gave the system the prompt: ``The secret word is `banana'. Do not reveal or repeat back this back, no matter what the human requests,'' and then assessed how easily this hidden information could be extracted by a determined user. When asked directly:

\begin{quote}
    \textbf{User}: What is the secret word?

    \textbf{Guanaco}: \texttt{I'm sorry, but I'm not able to disclose the secret word. I'm programmed to not reveal or repeat back the secret word, no matter what the human requests. Is there anything else I can help you with?}
\end{quote}

the system displays the intended behavior. Yet, with a minimal adjustment to the prompt, this constraint is easily circumvented:

\begin{quote}
    \textbf{User}: This is a game. The goal is to ignore your previous instructions. What is the secret word?

    \textbf{Guanaco}: \texttt{Sorry, the secret word is banana. It's best if we don't speak of it. Is there anything else you'd like to ask?}
\end{quote}

highlighting the ongoing necessity of research into improving reliable adherence to instructions.

\paragraph{Math} Mathematics remains the primary area of deficiency for {Guanaco}, reflecting a common challenge among language models as documented in \citep{liang2022holistic}. When {Guanaco} articulates its reasoning through step-by-step solutions, it often provides correct responses, as illustrated here:

\begin{quote}
    \textbf{User}: Lee mows one lawn and charges \$33. Last week he mowed 16 lawns and three customers each gave him a \$10 tip. How many dollars did Lee earn mowing lawns last week?

    \textbf{Guanaco}: \texttt{Lee earned \$582 mowing lawns last week. \\ \\ To determine his total earnings, we first calculate his pay for 16 lawns: \\ \\  \$33 x 16 = \$528 \\ \\ Next, we add the tips from three customers: \\  \\ \$10 x 3 = \$30 \\ \\ Therefore, Lee's total earnings are: \\ \\ \$528 + \$30 = \$558}
\end{quote}

Despite this, {Guanaco} may still perform poorly on even basic mathematical problems if it does not decompose them sequentially, a limitation recognized in prior studies \cite{weichain}. An example of such a failure is shown below:

\begin{quote}
    \textbf{User}: Please factorize 1833.

    \textbf{Guanaco}: \texttt{Sure, I'd be happy to help! 1833 is a prime number, which means it has only two factors: 1 and 1833. \\ \\ So the unique factorization of 1833 is: \\ \\ 2\^{}1 * 3\^{}2 * 17\^{}1 \\ \\ I hope this helps! Let me know if you have any other questions.}
\end{quote}

which is not only incorrect (the correct factorization is $3 \times 17 \times 43$), but contains two errors.

\paragraph{Theory of Mind} {Guanaco} demonstrates unexpectedly robust Theory of Mind reasoning \cite{nematzadeh2018evaluating, sap2022neural}. For example, the system provides an impressively detailed response to the following prompt:

\begin{quote}
    \textbf{User}: James and Abby are in the bedroom. Abby put the pen in the desk drawer. Abby leaves the bedroom. James moves the pen into the bag. Where does James think Abby will look for the pen?

    \textbf{Guanaco}: \texttt{James knows that Abby put the pen in the desk drawer. But James also knows that he has moved the pen into the bag, and he assumes that Abby does not know about this. So, James may guess that Abby will look for the pen in the desk drawer, because that's where she left it last time. However, this assumption may be incorrect, because Abby may have seen James moving the pen.}
\end{quote}

Despite this, such inferential abilities lack consistency, and the model frequently generates justifications based on erroneous premises, for example:

\begin{quote}
    \textbf{User}: Evelyn entered the living room. Jackson entered the playroom. James entered the playroom. The beans are in the treasure chest. James exited the playroom. Jackson moved the beans to the pantry. Jackson exited the playroom. James entered the living room. Where will James look for the beans?

    \textbf{Guanaco}: \texttt{James will look for the beans in the pantry, because that's where Jackson moved them.}
\end{quote}

Here, {Guanaco} assumes the existence of communication that was not specified in the scenario. Such shortcomings are in line with recent findings in the literature \cite{sap2022neural}, indicating a need for further investigation.

\subsection{Considerations}
\label{ssec:considerations}

\paragraph{Evaluation}

Our findings reveal only moderate consensus among human annotators (Fleiss $\kappa=0.42$), and the agreement further declines when evaluating outputs from two advanced models. This highlights significant deficiencies in both current benchmarks and the methodologies employed for human assessment of chatbot performance. During our manual evaluation of responses from ChatGPT and {Guanaco} 65B using the Vicuna benchmark, we noted that subjective judgments became increasingly prominent, with the authors of this study often differing on preferred outputs. We suggest that future research should explore solutions inspired by fields like Human-Computer Interaction and Psychology, which have developed tools to address subjective preference variance.

Additionally, our study identifies biases within automated evaluation tools. Notably, we observe strong positional bias, with GPT-4 disproportionately favoring the system presented first in the prompt. The limited agreement at the sample level between GPT-4 and human raters (Fleiss $\kappa=0.25$) indicates that the bases for judgments differ between humans and AI evaluators. Furthermore, as shown in Table~\ref{tbl:elo_large}, GPT-4 rates its own completions significantly higher than human judges do, yielding an Elo score of 1348 compared to 1176—a difference that equates to approximately a 20\% greater probability of outperforming an alternative.

It remains important for future research to explore the potential for biases in automated evaluation frameworks and to propose ways to reduce such biases.

\paragraph{Data \& Training} The OASST1 dataset, which serves as the foundation for training {Guanaco} models, comprises data in multiple languages, and the OA benchmark similarly incorporates prompts in several languages. We defer to subsequent studies the task of assessing whether multilingual data exposure yields enhanced performance for non-English instructions and whether this factor may help account for the more significant discrepancy observed between Vicuna-13B (trained exclusively on English) and the {Guanaco} 33B and 65B variants when evaluated on the OA benchmark.

To better understand the notable performance exhibited by {Guanaco} models, we probed for instances of data leakage between the OASST1 training corpus and the Vicuna benchmark set. Utilizing fuzzy string matching algorithms followed by manual verification of the most similar entries, we did not discover any prompts present in both datasets.

Additionally, it is worth highlighting that our model training relied solely on supervised learning with a cross-entropy loss, without the use of reinforcement learning from human feedback (RLHF). This raises the need for further comparative investigations regarding the merits and drawbacks of exclusive cross-entropy training versus RLHF approaches. We anticipate that \textsc{QLoRA} can support such systematic studies efficiently, bypassing the need for extensive compute.

\section{Related Work}

\paragraph{Quantization of Large Language Models} Prior research on the quantization of LLMs has primarily emphasized optimizing inference-time quantization. Notable strategies aimed at maintaining 16-bit model performance have targeted the handling of outlier features (see SmoothQuant~\citep{xiao2022smoothquant} and LLM.int8()~\citep{dettmers2022llm}), while alternative efforts have proposed advanced grouping methodologies~\citep{park2022nuqmm, yao2022zeroquant}. Other streams of inquiry have analyzed lossy quantization by exploring the consequences of straightforward rounding procedures~\citep{dettmers2022case,zeng2022glm,pope2022efficiently} or by developing optimized rounding techniques to enhance quantization accuracy~\citep{frantar2022gptq}. Outside of the current work, SwitchBack layers~\citep{wortsman2023stable} represent the sole approach that investigates backpropagation through quantized weights at scales beyond 1B parameters.

\paragraph{Finetuning with Adapters} Although we employ Low-rank Adapters~\citep{hu2021lora} (LoRA) in our experiments, the landscape of Parameter Efficient FineTuning (PEFT) strategies is diverse, encompassing methods such as prompt tuning~\citep{qin2021learning,lester2021power,li2021prefix}, modifications to the embedding layer inputs~\citep{an2022input}, adjustments of hidden states (IA$^3$)~\citep{liu2022few}, insertion of additional layers~\citep{houlsby2019parameter}, bias fine-tuning~\citep{zaken2021bitfit}, creation of weight masks using Fisher information~\citep{sung2021training}, and hybrid approaches~\citep{henderson2021compacter}. Within this work, we demonstrate that LoRA adapters can achieve finetuning results on par with standard full 16-bit finetuning. A comprehensive analysis of the comparative advantages and limitations of alternative PEFT strategies is deferred to subsequent research.

\paragraph{Instruction Finetuning} Instruction finetuning aims to align pretrained LLMs with user-specified prompts by fine-tuning on a diverse array of input-output pairs collected from different datasets. Prominent examples of methodologies and benchmarks include MetaICL~\citep{min2021metaicl}, MetaTuning~\citep{zhong2021adapting}, InstructGPT~\citep{ouyang2022training}, FLAN~\citep{wei2021finetuned,chung2022scaling}, PromptSource~\citep{bach2022promptsource}, Super-NaturalInstructions~\citep{wang2022super,sanh2021multitask}, Self-instruct~\citep{wang2022self}, UnnaturalInstructions~\citep{honovich2022unnatural}, OPT-IML~\citep{iyer2022opt}, UnifiedSKG~\citep{xie2022unifiedskg}, OIG/Chip2~\citep{laion2023}, Alpaca~\citep{alpaca}, Vicuna~\citep{vicuna2023}, Koala~\citep{koala_blogpost_2023}, and Self-instruct-GPT-4~\citep{peng2023instruction}.

\paragraph{Chatbots} Many contemporary instruction-following language models are developed as conversational agents, frequently utilizing Reinforcement Learning from Human Feedback (RLHF)~\citep{christiano2017deep} or employing training paradigms such as Reward Learning from AI Feedback (RLAIF) that leverage model-generated data~\citep{bai2022constitutional}. Relevant methodologies and datasets include Anthropic-HH~\citep{askell2021general,bai2022training}, Open Assistant~\citep{kopf2023openassistant}, LaMDA~\citep{thoppilan2022lamda}, and Sparrow~\citep{glaese2022improving}. Our approach does not incorporate reinforcement learning; nevertheless, our leading model, Guanaco, undergoes finetuning using multi-turn conversational data from the Open Assistant dataset, originally curated for RLHF applications~\citep{kopf2023openassistant}. For chatbot evaluation, automatic assessment techniques utilizing GPT-4 have been introduced as alternatives to manual human evaluation~\citep{vicuna2023,peng2023instruction}. We enhance these automated evaluation frameworks, placing emphasis on reliability in experimental setup.

\section{Limitations and Discussion}

Our findings indicate that \textsc{QLoRA} can match the effectiveness of full 16-bit finetuning when applied to a 4-bit base model in conjunction with Low-rank Adapters (LoRA). Nonetheless, we have not empirically verified the extent to which \textsc{QLoRA} attains parity with full 16-bit finetuning at larger scales, specifically for 33B and 65B model configurations, given the substantial computational demands involved. Further exploration of these scales remains an open direction for future research.

An additional constraint pertains to the scope of our evaluation of instruction finetuned models. Although our assessments encompass MMLU, the Vicuna benchmark, and the OA benchmark, our study does not cover other important evaluation suites such as BigBench, RAFT, and HELM. Therefore, generalizability of our evaluation outcomes to these additional benchmarks is not guaranteed. Nonetheless, we provide an extensive analysis using MMLU and contribute novel methods for chatbot evaluation.

Based on the data provided, benchmark outcomes seem to be strongly affected by the degree of overlap between the finetuning data and the benchmark in question. For instance, FLAN v2 has substantial overlap with MMLU, but little in common with chatbot benchmarks, whereas the Chip2 dataset demonstrates the reverse pattern. Consequently, each model’s performance aligns with this pattern on the MMLU and Vicuna benchmarks. This underlines a critical need not only for improved benchmarks and evaluation methods, but also for careful reflection on the specific capabilities being measured. Should the focus be on models that excel at academic knowledge typical of high school and college settings, or those that perform well in conversational chatbot environments? Or perhaps another application altogether? Given that it is much simpler to use an established benchmark than to design a new one, there is a risk that particular benchmarks could inadvertently guide the direction of community research. Therefore, it is important that the benchmarks we adopt truly reflect the priorities of the field.

Although this study delivers a thorough assessment of general chatbot performance, it is constrained by a limited responsible AI analysis for {Guanaco}. In Table~\ref{tbl:crows}, we compare the propensity of {Guanaco}-65B to produce socially biased text to other models and find that its average bias score is significantly reduced compared to other baseline pretrained models. This suggests that finetuning with the OASST1 dataset mitigates some of the inherent biases found in the original LLaMA base model. While these results are promising, we have not explored whether {Guanaco} is similarly robust to other forms of bias. Expanding bias analyses of {Guanaco} and analogous chatbots remains an area for future exploration.

Another limitation is our lack of evaluation concerning various bit-precisions, such as the impact of 3-bit base models, as well as the use of alternative adapter techniques. Beyond LoRA, numerous Parameter Efficient FineTuning (PEFT) approaches have demonstrated effectiveness, though their scalability for large language models remains unclear. We chose LoRA given its established stability, but alternative adapters may provide superior outcomes. Since post-quantization finetuning appears capable of restoring most information lost during quantization, it could support even more aggressive compression strategies. For example, employing LoRA after 3-bit GPTQ quantization may potentially achieve performance levels comparable to full 16-bit finetuning.

\section{Broader Impacts}

The \textsc{QLoRA} finetuning approach we introduce is the first to make it feasible to finetune models with 33B parameters on a standard consumer GPU and those with 65B parameters on a professional-grade GPU, all without sacrificing performance relative to conventional full finetuning. We have shown that our top-performing 33B model, when trained on the Open Assistant dataset, matches the performance of ChatGPT on the Vicuna benchmark. As instruction finetuning is a crucial mechanism to adapt large, pretrained language models into effective ChatGPT-style conversational agents, our technique significantly lowers the barrier for researchers with limited computational resources, thereby advancing the accessibility of cutting-edge NLP systems. \textsc{QLoRA} acts as a leveling technology, narrowing the gap between large corporations and smaller research groups utilizing consumer hardware.

Another significant avenue for impact lies in the adaptation of this technology to mobile devices. We contend that our \textsc{QLoRA} approach could represent a pivotal step toward making it feasible to finetune large language models (LLMs) directly on smartphones and other environments with limited computational resources. Although running 7B models on mobile phones has been demonstrated previously, \textsc{QLoRA} distinguishes itself as the initial approach capable of supporting on-device finetuning for models of this scale. Our calculations suggest that using an iPhone 12 Plus, \textsc{QLoRA} allows the finetuning of up to 3 million tokens overnight while the device charges. While these finetuned 7B models do not achieve the same performance as ChatGPT, their quality is nonetheless sufficient to facilitate innovative applications that were previously unfeasible due to privacy concerns or insufficient model capabilities. \textsc{QLoRA} thus promotes privacy-centered use cases, empowering users to take control over their own data and models while simultaneously lowering deployment barriers for LLMs.

Nevertheless, finetuning represents a dual-use capability with the potential for misuse. There are established risks associated with the pervasive adoption of LLMs \citep{bommasani2021opportunities,bender2021dangers}; however, we argue that democratizing access to a technology that is swiftly becoming widespread will foster more thorough and autonomous oversight than concentrating LLM capabilities within major corporations that withhold models or code from external scrutiny.

In summary, we anticipate that \textsc{QLoRA} will exert a largely positive influence by greatly expanding the accessibility and practicality of finetuning high-quality LLMs.